\documentclass[../../../include/open-logic-section]{subfiles}

\begin{document}

\olfileid{his}{set}{mythology}

\olsection{اسطوره‌ای بیش از تاریخ؟}

اکنون که با بیش از یک سده فاصله به این رویدادها می‌نگریم، باید مراقب
باشیم این داوری‌ها را بی‌چون‌وچرا نپذیریم. نتایج بی‌تردید تازه، هیجان‌انگیز
و شگفت‌آور بودند؛ اما واقعاً تا چه اندازه تکان‌دهنده بودند؟ و آیا به‌راستی
نشان می‌دادند که نباید بر شهود هندسی تکیه کنیم؟

دربارهٔ میزان شگفتی، \citet{Gouvea2011} خاطرنشان می‌کند که عبارت مشهور
کانتور در نامه به ددکیند، «\emph{je le vois, mais je ne le crois pas}»،
تا حد زیادی بیرون از بافت خود نقل می‌شود. بخش بیشتری از آن بافت را در زیر
می‌آوریم (به نقل از \citeauthor{Gouvea2011}):
\begin{quote}
خواهش می‌کنم شور مرا نسبت به این موضوع ببخشید، اگر خواسته‌های فراوانی از
مهربانی و شکیبایی شما دارم؛ مطالبی که به‌تازگی برایتان فرستادم حتی برای
خودم چنان نامنتظر و تازه‌اند که تا از شما، دوست ارجمند، حکمی دربارهٔ
درستی‌شان نگیرم آرامش نخواهم یافت. تا هنگامی که با من موافقت نکرده‌اید،
تنها می‌توانم بگویم: \emph{je le vois, mais je ne le crois pas.}
\end{quote}
کانتور می‌دانست نتیجه‌اش «چنان نامنتظر و تازه» است؛ اما تردید داریم که
هرگز آن را \emph{باورنکردنی} دانسته باشد. چنان‌که
\citeauthor{Gouvea2011} می‌گوید، او صرفاً از ددکیند می‌خواست برهانی را که
پیشنهاد کرده بود بررسی کند.

اما دربارهٔ شهود هندسی: پئانو منحنی پُرکنندهٔ فضای خود را بدون هیچ شکلی
منتشر کرد. ولی هیلبرت هنگام انتشار منحنی‌اش مقصود خود را توضیح داد: او
راهی روشن برای فهم نتیجهٔ پئانو به خوانندگان عرضه می‌کرد، به شرط آنکه
«از شهود هندسی زیر یاری بگیرند»؛ سپس رشته‌ای از \emph{شکل‌ها} درست مانند
شکل‌های \olref[his][set][pathology]{sec} آورد.

به‌طور کلی‌تر، گرچه شکل‌ها در برهان‌های منتشرشده تا حدی از رواج افتاده‌اند،
نمی‌توان این واقعیت را نادیده گرفت که ریاضی‌دانان هنگام اثبات مطالب
\emph{اغلب} از شکل استفاده می‌کنند. (به‌اختصار: ریاضی‌دان خوب می‌داند چه
وقت می‌تواند بر شهود هندسی تکیه کند.)

خلاصه آنکه تبلیغات را باور نکنید؛ یا دست‌کم بی‌چون‌وچرا نپذیریدشان. برای
مطالعهٔ بیشتر می‌توانید \citet{Giaquinto2007} را بخوانید.

\end{document}
